\documentclass[12pt, a4paper, oneside]{ctexart}
\usepackage{amsmath, amsthm, amssymb, bm, graphicx, hyperref, mathrsfs}

\title{\textbf{科研课堂课程作业}}
\author{黄伟宁}
\date{\today}
\linespread{1.5}
\newcounter{problemname}
\newenvironment{problem}{\stepcounter{problemname}\par\noindent\textbf{题目\arabic{problemname}. }}{\\\par}
\newenvironment{solution}{\par\noindent\textbf{解答. }}{\\\par}
\newenvironment{note}{\par\noindent\textbf{题目\arabic{problemname}的注记. }}{\\\par}
\usepackage{geometry}
\geometry{left=2.54cm,right=2.54cm,top=3.18cm,bottom=3.18cm}

\begin{document}

\maketitle

\begin{problem}
验证KdV方程$ u_{t} +6uu_{x} +u_{xxx} =0 $ 的Lax对.
\end{problem}

\begin{solution}
    考虑如下两个线性算子$ L,A $
    $$
        \begin{cases}
            L=\frac{\partial ^{2} }{\partial x^{2} }+u \\
            A=-4\frac{\partial ^{3} }{\partial x^{3} }-6u\frac{\partial }{\partial x}-3u_{x}
        \end{cases}
    $$
    那么我们有
    $$
        \begin{cases}
            L\varphi=\lambda\varphi \\
            \varphi_{t} =A\varphi
        \end{cases}
    $$
    不难发现
    \begin{align*}
         & L_{t} \varphi+L\varphi_{t} =\lambda\varphi_{t} =\lambda A\varphi \\
         & \Rightarrow L_{t} \varphi=(AL-LA)\varphi
    \end{align*}
    即Lax对满足
    $$
        L_{t} =AL-LA\equiv[A,L]
    $$
    我们有
    \begin{align*}
         & L_{t} =u_{t}                                                                                                                                                                                                                      \\
         & AL=-(4\frac{\partial ^{5} }{\partial x^{5} }+10u\frac{\partial ^{3} }{\partial x^{3} }+15u_{x} \frac{\partial ^{2} }{\partial x^{2} }+4u_{xxx}+12u_{xx} \frac{\partial }{\partial x}+6u^{2} \frac{\partial }{\partial x}+9uu_{x}) \\
         & LA=-(4\frac{\partial ^{5} }{\partial x^{5} }+12u_{xx} \frac{\partial }{\partial x}+15u_{x} \frac{\partial ^{2} }{\partial x^{2} }+10u\frac{\partial ^{3} }{\partial x^{3} }+6u^{2} \frac{\partial }{\partial x}+3u_{xxx}+3uu_{x}) \\
         & AL-LA=-u_{xxx} -6uu_{x}
    \end{align*}
    从而有
    $$
        u_{t} +6uu_{x} +u_{xxx} =0
    $$
    此即所给KdV方程.
\end{solution}

\begin{problem}
验证KdV方程$ u_{t} +6uu_{x} +u_{xxx} =0 $ 的DT.
\end{problem}

\begin{solution}
    略.
\end{solution}

\end{document}